% =============================================================================
%  MỞ ĐẦU (3-5 trang)
% =============================================================================
\chapter*{MỞ ĐẦU}
\addcontentsline{toc}{chapter}{MỞ ĐẦU}
\markboth{MỞ ĐẦU}{}

\section*{1. Tính cấp thiết của đề tài}
\addcontentsline{toc}{section}{1. Tính cấp thiết của đề tài}

Mạng thông tin di động thế hệ thứ sáu (6G) được kỳ vọng sẽ vượt xa các thế
hệ trước về thông lượng (lên tới $1$~Tbps), độ trễ siêu thấp (dưới
$0{,}1$~ms) và mật độ kết nối khổng lồ ($10^7$ thiết bị/km$^2$)
\cite{saad2020vision6g}, \cite{tataria20216g}. Để đáp ứng các yêu cầu khắt khe đó,
nhiều công nghệ then chốt được nghiên cứu, trong đó hai hướng nổi bật là
Đa truy nhập theo phân tách tốc độ (Rate-Splitting Multiple Access
RSMA) và Bề mặt phản xạ thông minh có khả năng truyền-phản xạ đồng thời
(Simultaneously Transmitting And Reflecting Reconfigurable Intelligent
Surface, STAR-RIS) \cite{mao2022rsma}, \cite{liu2021starris}, \cite{mu2022simultaneously}.

RSMA cho phép trạm gốc phân tách thông điệp của mỗi người dùng thành một
phần chung (common) và một phần riêng (private), sau đó truyền
đồng thời qua cùng tài nguyên thời gian-tần số. Tại phía thu, người dùng
giải mã luồng chung trước, áp dụng khử nhiễu liên tiếp (SIC) và sau đó giải
mã luồng riêng. Cách thiết kế này tạo nên khả năng quản lý nhiễu mềm dẻo
hơn so với các kỹ thuật trực giao (OMA) và không trực giao truyền thống
(NOMA), đặc biệt là trong điều kiện thông tin trạng thái kênh không hoàn
hảo \cite{mao2018rsma}, \cite{mao2022rsma}.

STAR-RIS là sự mở rộng của RIS truyền thống, gồm $N$ phần tử thụ động có thể
đồng thời điều chỉnh pha cả tín hiệu phản xạ và tín hiệu truyền qua
mặt phẳng RIS. Đặc tính này biến vùng phủ sóng từ nửa không gian
$180^\circ$ thành toàn không gian $360^\circ$, mở ra khả năng phục vụ
người dùng ở hai phía của bề mặt thông minh \cite{liu2021starris}, \cite{mu2022simultaneously}.

Tuy nhiên, khi tích hợp RSMA và STAR-RIS, bài toán phân bổ tài nguyên
trở nên cực kỳ phức tạp: ta phải đồng thời tối ưu hóa (i) công suất phát
giữa luồng chung và các luồng riêng, (ii) tỷ lệ phân chia luồng chung cho
từng người dùng, (iii) biên độ và pha của mỗi phần tử STAR-RIS ở cả hai
chế độ phản xạ và truyền qua. Bài toán này có tính phi lồi cao do các
hàm SINR phụ phi tuyến vào biến công suất và vào tích vô hướng phức của
vector pha. Các phương pháp tối ưu truyền thống như hạ dần khối tọa độ
(BCD), nới lỏng nửa xác định (SDR) hoặc tối ưu xen kẽ (AO) tuy có thể
cho lời giải gần tối ưu nhưng đòi hỏi chi phí tính toán lớn, không phù
hợp cho triển khai thời gian thực trong môi trường kênh biến thiên nhanh
\cite{wu2020smart}.

Học tăng cường sâu (DRL), và đặc biệt là phiên bản đa tác tử MADDPG
cung cấp một con đường khác: thay vì giải tường minh bài toán tối ưu, tác tử
học chính sách phân bổ tài nguyên trực tiếp từ tương tác với môi trường
\cite{sutton2018reinforcement}, \cite{lillicrap2016continuous}, \cite{lowe2017multi}.
Sau khi huấn luyện, suy luận chỉ cần một lần chạy tiến (forward pass) của
mạng nơ-ron, có độ phức tạp gần như hằng số, lý tưởng cho yêu cầu thời
gian thực của 6G \cite{noman2026drl6g}, \cite{hieu2021optimal}. Vì các lý do trên,
việc nghiên cứu ứng dụng MADDPG để tối ưu phân bổ tài nguyên trong
mạng STAR-RIS hỗ trợ RSMA mang ý nghĩa khoa học và thực tiễn rõ rệt.

\section*{2. Mục tiêu nghiên cứu}
\addcontentsline{toc}{section}{2. Mục tiêu nghiên cứu}

\noindent\textbf{Mục tiêu tổng quát:} Xây dựng và đánh giá một khuôn khổ
học tăng cường sâu đa tác tử cho bài toán tối ưu phân bổ tài nguyên đường
xuống trong mạng STAR-RIS hỗ trợ RSMA, sao cho tối đa hóa tổng tốc độ hệ
thống đồng thời tối ưu cân bằng giữa thông lượng và mức đáp ứng chất lượng
dịch vụ (QoS) của từng người dùng thông qua một hàm mục tiêu thay thế
(surrogate) dạng phạt/Lagrangian của ràng buộc QoS, tương ứng với ràng buộc
kỳ vọng $\mathbb{E}[R_k] \geq R_{\min}$.

\noindent\textbf{Mục tiêu cụ thể:}
\begin{enumerate}[label=(\arabic*),leftmargin=*]
\item Xây dựng mô hình hệ thống đầy đủ: kênh truyền (trực tiếp, gián tiếp
qua STAR-RIS, suy hao đường truyền log-distance, pha-đinh Rayleigh),
mô hình tín hiệu RSMA, và bài toán tối ưu phân bổ tài nguyên (P0) với
đầy đủ các ràng buộc về công suất, bảo toàn năng lượng STAR-RIS, phân
chia luồng chung và QoS.
\item Đề xuất khuôn khổ MADDPG với ba tác tử chuyên dụng (BS Power, STAR-RIS
Reflection, STAR-RIS Transmission) cùng hai cải tiến: (i) tham số hóa
pha residual dựa trên nghiệm phân tích bậc một và (ii) hàm phần thưởng
dạng Lagrangian với các hệ số phạt $\lambda_k$ riêng cho từng người dùng,
được cập nhật bằng phương pháp gradient đối ngẫu có phép chiếu
(projected dual gradient) trên khe hở ràng buộc có dấu
$c_k = R_{\min} - R_k$, để cân bằng tốc độ tổng và QoS.
\item Cài đặt mô phỏng và đánh giá toàn diện hiệu năng so với các baseline
(DDPG đơn tác tử, TD3, TD3 cân bằng tham số, PPO, AO-Grid, Random-RIS,
Fixed-RIS, NoRIS) trên các phương diện: hội tụ huấn luyện, tổng tốc độ,
mức đáp ứng QoS, độ trễ suy luận và phân tích loại trừ (ablation) đối
với chế độ STAR-RIS.
\end{enumerate}

\section*{3. Đối tượng nghiên cứu}
\addcontentsline{toc}{section}{3. Đối tượng nghiên cứu}

\begin{itemize}[leftmargin=*]
\item Mạng vô tuyến đường xuống STAR-RIS hỗ trợ RSMA, gồm trạm gốc, bề mặt
phản xạ thông minh và các thiết bị người dùng ở cả hai phía của bề mặt.
\item Bài toán tối ưu phân bổ đa biến (công suất, biên độ, pha, tỷ lệ tách
luồng) cho mạng nêu trên.
\item Thuật toán học tăng cường sâu MADDPG và các thuật toán DRL tham
chiếu (DDPG, TD3, PPO).
\item Các phương pháp tối ưu cổ điển làm mốc so sánh: heuristic tối ưu xen
kẽ dạng lưới thô (AO-Grid), tìm kiếm cục bộ AO lai (Hybrid AO Local
Search, SLSQP + projected gradient) và nghiệm phân tích bậc một
(MaxMin-Aligned). Các phương pháp này là mốc tham chiếu cục bộ, không
phải cận trên tối ưu.
\end{itemize}

\section*{4. Phạm vi nghiên cứu}
\addcontentsline{toc}{section}{4. Phạm vi nghiên cứu}

\noindent\textbf{Phạm vi nội dung:} Tập trung vào bài toán tối ưu sum-rate
với ràng buộc QoS cho cấu hình truyền đường xuống MISO với trạm gốc
$M=4$ anten, một STAR-RIS có $N=32$ phần tử ở chế độ chia sẻ năng lượng
(ES), $K=4$ người dùng (gồm 3 người dùng vùng phản xạ và 1 người dùng
vùng truyền qua chịu suy hao chướng ngại $-25$~dB). Không xem xét chế độ
chuyển đổi MS/TS, RIS chủ động, RIS đa lớp hay các kịch bản multi-cell.

\noindent\textbf{Phạm vi không gian:} Mô phỏng số học bằng phần mềm
Python/PyTorch trên máy tính cá nhân, không triển khai phần cứng vật lý.

\noindent\textbf{Phạm vi thời gian:} Đề tài được nghiên cứu, mô phỏng và
hoàn thiện trong giai đoạn 2025--2026.

\section*{5. Phương pháp nghiên cứu}
\addcontentsline{toc}{section}{5. Phương pháp nghiên cứu}

\begin{itemize}[leftmargin=*]
\item \textbf{Phương pháp mô hình hóa toán học:} Xây dựng mô hình hệ thống,
phát biểu bài toán tối ưu (P0) và ánh xạ sang quá trình ra quyết định
Markov (MDP).
\item \textbf{Phương pháp học máy:} Thiết kế kiến trúc mạng nơ-ron Actor và
Critic, áp dụng cơ chế huấn luyện tập trung, thực thi phân tán (CTDE)
trong MADDPG, sử dụng các kỹ thuật nâng cao như mạng mục tiêu (target
networks) với cập nhật mềm, nhiễu khám phá Ornstein--Uhlenbeck (OU),
chuẩn hóa lớp (LayerNorm) và cắt gradient.
\item \textbf{Phương pháp thực nghiệm:} Mô phỏng đa hạt giống huấn luyện
(8 hạt giống độc lập) với khoảng tin cậy $95\%$ Student-t tính qua các
hạt giống huấn luyện; đánh giá trên tập kịch bản kiểm tra khóa
(ScenarioBank) tách biệt khỏi tập validation dùng cho tinh chỉnh; so
sánh với các baseline DDPG, TD3, PPO, AO-Grid, Random, Fixed, NoRIS;
phân tích loại trừ và kiểm định ý nghĩa thống kê theo cặp (paired
t-test, permutation test) với hiệu chỉnh Holm--Bonferroni cho so sánh
bội.
\end{itemize}

\section*{6. Ý nghĩa khoa học và thực tiễn}
\addcontentsline{toc}{section}{6. Ý nghĩa khoa học và thực tiễn}

\noindent\textbf{Về mặt lý luận:}
\begin{itemize}[leftmargin=*]
\item Mở rộng bài toán P0 cổ điển bằng việc bổ sung ràng buộc QoS từng
người dùng và xem xét biên độ STAR-RIS như biến tối ưu thay vì cố định.
\item Đề xuất tham số hóa pha residual, một kỹ thuật giúp giảm
không gian khám phá của tác tử pha và cải thiện đáng kể tốc độ hội tụ.
\item Đề xuất hàm phần thưởng dạng Lagrangian surrogate với các hệ số
$\lambda_k$ riêng cho từng người dùng, được cập nhật bằng projected dual
gradient trên khe hở ràng buộc có dấu, cho phép tự điều chỉnh cân bằng
giữa tổng tốc độ và mức đáp ứng QoS mà không cần điều chỉnh thủ công.
\item Đề xuất lịch trình huấn luyện hai pha đóng băng hệ số phạt
(two-stage dual freeze), một heuristic làm cho hàm phần thưởng trở nên
dừng ở giai đoạn cuối huấn luyện để đường cong huấn luyện dễ diễn giải;
lưu ý rằng vùng phẳng (plateau) của phần thưởng không tự nó chứng minh
chính sách đã hội tụ.
\end{itemize}

\noindent\textbf{Về mặt thực tiễn:}
\begin{itemize}[leftmargin=*]
\item Kết quả mô phỏng cho thấy tiềm năng triển khai thực tế của MADDPG
cho phân bổ tài nguyên thời gian thực trong 6G: sau huấn luyện, mỗi
quyết định chỉ cần một lần chạy tiến của mạng nơ-ron (độ trễ đo được
$\approx 1{,}3$~ms/hành động ở pipeline cũ, thấp hơn một bậc so với
heuristic AO-Grid giải lại cho từng thể hiện kênh; số liệu chính thức
sẽ dùng cặp bảng CPU/GPU đo riêng của pipeline mới), đồng thời định
lượng lợi thế mở rộng của kiến trúc đa tác tử khi số phần tử STAR-RIS
tăng tới $N=128$.
\item Cung cấp một thư viện mã nguồn (Python/PyTorch) sẵn sàng tái sử dụng
cho các nghiên cứu tiếp theo về phân bổ tài nguyên trong mạng RIS hoặc
RSMA.
\item Cấu hình mô phỏng (STAR-RIS $N=32$, $K=4$, $P_{\max}=30$~dBm) là một
chuẩn so sánh tham chiếu hữu ích cho cộng đồng nghiên cứu trong lĩnh
vực này.
\end{itemize}

\section*{7. Kết cấu của Luận văn}
\addcontentsline{toc}{section}{7. Kết cấu của Luận văn}

Ngoài phần Mở đầu và Kết luận, Luận văn gồm bốn chương như sau:
\begin{itemize}[leftmargin=*]
\item \textbf{Chương 1, Cơ sở lý thuyết:} Trình bày tổng quan về mạng 6G,
nguyên lý RSMA, công nghệ STAR-RIS, khuôn khổ DRL và thuật toán MADDPG.
\item \textbf{Chương 2, Mô hình hệ thống và bài toán tối ưu:} Xây dựng
mô hình kênh truyền, mô hình tín hiệu RSMA, phát biểu bài toán tối ưu
(P0) và ánh xạ sang MDP.
\item \textbf{Chương 3, Phương pháp đề xuất MADDPG:} Trình bày chi tiết
kiến trúc ba tác tử chuyên dụng, mạng nơ-ron Actor-Critic, các cải
tiến (residual phase, Lagrangian reward thích nghi) và quy trình huấn
luyện.
\item \textbf{Chương 4, Kết quả mô phỏng và đánh giá:} Trình bày toàn bộ
kết quả mô phỏng số học, so sánh với các baseline, phân tích loại trừ
và thảo luận về ý nghĩa thực tiễn.
\end{itemize}
